\section{Conclusion}
\label{sec:conclusion}

We introduced a dual-constraint bisection framework for partisan-proportional
redistricting. It extends minimum-edge-cut algorithms with two vertex weights:
population and Democratic votes.
The central contribution is a closed-form formula~\eqref{eq:tpwgts-formula}
for the top-level partition targets that express a proportional seat allocation:
\[
  \mathtt{tpwgts} = \left[\frac{k_D}{k},\; 1 - \frac{k_R}{2dk},\;
                          \frac{k_R}{k},\; \frac{k_R}{2dk}\right]
\]

The partisan Lorenz curve gives a non-contiguous feasibility diagnostic for
this target.
The partisan signal strength $\sigma$ measures how much the neutral geographic
distribution must be ``pulled'' toward the proportional target --- from
$\sigma \approx 0$ (the proportional top-level target is almost neutral) to
$\sigma > 0.2$ (the target requires substantial partisan concentration).

\paragraph{Empirical tradeoff results.}
Running the T.5 constraint at $\eta \in \{1.05, 1.10, 1.20\}$ across 6 competitive
states produces non-monotone, state-specific results:
Wisconsin achieves $-0.3$\,pp at $\eta = 1.10$ (near-exact proportionality, from
$-12.8$\,pp baseline), but $-25.3$\,pp at $\eta = 1.20$ (worse than baseline).
North Carolina worsens at all $\eta$ values ($-6.5$\,pp baseline $\to$ $-13.6$\,pp):
the HH 7:7 ratio forced by the constraint underperforms T.2's natural 6:8 ratio.
Arizona shows no change at any $\eta$ (geography insensitive).
There is no universal optimal $\eta$.

\paragraph{The core finding: the proportionality paradox.}
Competitive states (46--55\,\% Democratic) have $\sigma < 0.05$ --- the
proportional tpwgts target is essentially the same as the neutral equal-distribution
target. Georgia needs to shift 0.1\,\% of D votes from neutral; Wisconsin needs
0.3\,\%. Yet both have 12--14\,pp gaps under geography-only algorithms.

By Theorem~\ref{thm:invariance} (Tradeoff invariance), varying the partisan
constraint from loose ($\eta = \infty$) to strict ($\eta = 1.0$) cannot close
more than $C \cdot \sigma$ of the gap for competitive states. For Georgia
($\sigma = 0.001$) and Wisconsin ($\sigma = 0.003$), the target shift itself is
tiny compared with the observed 12--14\,pp gaps.

\textbf{The proportionality problem in competitive states is a Lorenz feasibility
problem, not a target calibration problem.}
The top-level target is easy to write down; the hard part is finding contiguous
downstream districts that realise it.

The states that most need proportional redistricting are the states where
proportionality is most expensive.
Extreme Democratic geographic concentration --- Milwaukee, Madison, Atlanta,
Boston --- produces large Rodden gaps \emph{and} makes proportional packing
geometrically difficult.
There is no geographic algorithm that can close the proportionality gap
in these states without explicitly incorporating partisan data.
This is a fact about geography.

\paragraph{What this means.}
For states with low $\sigma$ (cheap proportionality): geography-only algorithms
are nearly proportional. Small improvements to the geographic constraint
(T.2-style area balance or T.1-style isoperimetric normalisation) may close
the remaining gap without partisan data.

For states with high $\sigma$ (expensive proportionality): the proportionality
gap is a structural feature of the state's population distribution.
A court reviewing such a state's redistricting for partisan fairness should
understand that even a perfectly neutral algorithm produces a large gap ---
the gap is not evidence of manipulation, but of geography.

\paragraph{Future work.}
Three directions merit further investigation:
\begin{enumerate}
\item \textbf{Recursive proportionality}: applying the T.5 constraint at
  every level of the bisection tree and recording child-level feasibility
  witnesses, rather than using the current top-split-only implementation.

\item \textbf{Continuous tolerance curves}: empirical measurement of how the
  proportionality gap decreases as $\mathtt{ubvec}[1]$ is tightened, for each
  state in the 50-state sweep.

\item \textbf{Multi-member districts}: for expensive-proportionality states,
  characterising how many seats per district are needed before geographic
  constraints stop blocking proportionality --- connecting this analysis to
  the Huntington-Hill minimum from T.4.
\end{enumerate}
